\begin{abstract}
Audio-visual (AV) deepfake detectors increasingly build on frozen foundation encoders, yet their robustness is almost always reported empirically, without per-sample guarantees. We present CertAV, a certification framework that applies randomized smoothing in the joint feature space of frozen DINOv2 and Whisper encoders. On FakeAVCeleb, CertAV attains $92.5\%$ clean accuracy and $88.0\%$ certified accuracy at $\ell_2$ radius $1.0$, with under $1\%$ abstention. We then explain \emph{why} this works: under a subspace-invariance property of the detector head, the smoothed classifier's certified radius is governed entirely by the on-manifold component of the noise, whose intrinsic dimension $d\!\approx\!80$ is far below the ambient dimension $D\!=\!768$. This yields a robustness amplification factor $\sqrt{D/d}$ (Theorem~1), a principled encoder-selection criterion (Corollary~1), and a prediction that the certifiability is set by representation geometry rather than by training --- confirmed by a no-noise baseline that certifies within $1.9\%$ of the noise-trained model. We validate the scaling across five encoder families and exploit the geometry with manifold-aware anisotropic smoothing, which at equalized noise budget achieves a $3.4\times$ larger on-manifold certified radius ($7.63$ vs.\ $2.22$) with zero abstention. Finally, robust conformal prediction supplies a distributional coverage guarantee for every sample, including those CertAV abstains on. We release \texttt{av-robustbench} to make the full pipeline reproducible.
\end{abstract}
